\documentclass[conference]{IEEEtran}
\IEEEoverridecommandlockouts
% The preceding line is only needed to identify funding in the first footnote. If that is unneeded, please comment it out.
\usepackage{cite}
\usepackage{amsmath,amssymb,amsfonts}
\usepackage{algorithmic}
\usepackage{graphicx}
\usepackage{textcomp}
\usepackage{xcolor}
\def\BibTeX{{\rm B\kern-.05em{\sc i\kern-.025em b}\kern-.08em
    T\kern-.1667em\lower.7ex\hbox{E}\kern-.125emX}}
\begin{document}

\title{Comparing Critic and User Movie Reviews Using Data Mining}


\author{\IEEEauthorblockN{1\textsuperscript{st} Brady Lauritsen}
  \IEEEauthorblockA{\textit{University of Nebraska Lincoln} \\
    \textit{Computer Science | CSCE 474}\\
    Lincoln, NE, USA} \\

  \and
  \IEEEauthorblockN{2\textsuperscript{nd} Sagun Karki}
  \IEEEauthorblockA{\textit{University of Nebraska Lincoln} \\
    \textit{Computer Science \& Data Science}\\
    Lincoln, NE, USA} \\

  \and
  \IEEEauthorblockN{3\textsuperscript{rd} Given Name Surname}
  \IEEEauthorblockA{\textit{dept. name of organization (of Aff.)} \\
    \textit{name of organization (of Aff.)}\\
    City, Country} \\

  \and
  \IEEEauthorblockN{4\textsuperscript{th} Given Name Surname}
  \IEEEauthorblockA{\textit{dept. name of organization (of Aff.)} \\
    \textit{name of organization (of Aff.)}\\
    City, Country} \\
}

\maketitle

\begin{abstract}
  Movie ratings from professional critics and general audiences often differ significantly.
  This project explores the relationship between critic reviews from Rotten Tomatoes and
  user reviews from Letterboxd using data mining techniques. The goal is to identify patterns,
  analyze sentiment differences, and determine whether critic ratings can predict user ratings.
  By combining multiple datasets and performing preprocessing, correlation analysis, and
  predictive modeling, we aim to better understand how critics and audiences evaluate films.
\end{abstract}

\begin{IEEEkeywords}
  data mining, sentiment analysis, movie reviews, correlation analysis, prediction
\end{IEEEkeywords}

\section{Introduction}
Briefly introduce the project topic, why it is important, and what problem you are trying to solve.

\section{Objectives}
List the goals of the project.

\begin{itemize}
  \item Objective 1:
  \item Objective 2:
  \item Objective 3:
  \item Objective 4:
\end{itemize}

\section{Datasets}
Describe the datasets that will be used in this project.

\subsection{Dataset 1}
\begin{itemize}
  \item Source: Kaggle \cite{leone2020rotten}
  \item Number of records: 17k movies
  \item Description: Rotten Tomatoes dataset.
\end{itemize}

\subsection{Dataset 2}
\begin{itemize}
  \item Source:
  \item Number of records:
  \item Description:
\end{itemize}

\section{Data Processing}
Explain how the data will be prepared before analysis.

\begin{itemize}
  \item Data cleaning
  \item Handling missing values
  \item Data integration
  \item Feature selection
\end{itemize}

\section{Methods}
Describe the data mining or machine learning techniques that will be used.



\section{Evaluation}
Explain how the results of the project will be evaluated.

\begin{itemize}
  \item Accuracy
  \item Correlation metrics
  \item Visualization
\end{itemize}

\section{Implementation Plan}
Describe the tools and timeline for the project.

\subsection{Tools}


\subsection{Timeline}

\section{Team Contributions}
Briefly describe what each team member is responsible for.

\begin{itemize}
  \item Brady Lauritsen:
  \item Sagun Karki:
  \item Team Member 3:
  \item Team Member 4:
\end{itemize}

\section{Conclusion}
Summarize what the project aims to accomplish.

\bibliographystyle{IEEEtran}
\bibliography{references}
\end{document}

\end{document}